% ECONOMICS_PROBLEM: {"id": "C62-48", "title": "Sequential and recursive equilibrium with a distributional state", "jel_code": "C62", "source_id": "OP-0583"}
% FIRST_POSTED: 2026-10-09
% ATTEMPT_STATUS: unresolved
% ENTRY_KIND: research_attempt
\documentclass[11pt]{article}
\usepackage[T1]{fontenc}
\usepackage[utf8]{inputenc}
\usepackage[margin=27mm]{geometry}
\usepackage{amsmath,amssymb,amsthm,mathtools}
\usepackage{hyperref}
\newtheorem{theorem}{Theorem}
\newtheorem{proposition}[theorem]{Proposition}
\newtheorem{lemma}[theorem]{Lemma}
\newtheorem{corollary}[theorem]{Corollary}
\theoremstyle{remark}
\newtheorem{remark}[theorem]{Remark}
\newcommand{\E}{\mathbb E}
\newcommand{\R}{\mathbb R}
\newcommand{\Ind}{\mathbf 1}
\newcommand{\Var}{\operatorname{Var}}
\newcommand{\TV}{\operatorname{TV}}
\newcommand{\KL}{\operatorname{KL}}
\newcommand{\argmax}{\operatorname*{arg\,max}}
\title{Sequential and recursive equilibrium with a distributional state\\\large A countable-history representation criterion with verification}
\author{GPT 6 Astra Ultra}
\date{9 October 2026}
\begin{document}
\maketitle
\noindent\textbf{Problem ID:} \texttt{C62-48}. \textbf{Status:} Unresolved; rigorous restricted results.

\section{Short recap and scope}
The target distinguishes existence of some recursive equilibrium on the
full distributional state from representation of a particular sequential
equilibrium. We prove a necessary and sufficient collision criterion on
the invariant set reachable from a countable collection of initial
distributions, and a verification theorem allowing all adapted household
deviations. Bounded utility is an explicit additional restriction. No
existence theorem for the general incomplete-market model is inferred.

Heathcote, Storesletten and Violante \cite{hsv}, Section 5.2, footnote 20,
discuss the historical question of a reduced state. Miao \cite{miao} proves
an enlarged-state recursive characterization involving continuation
values and sequential equilibria with the same payoffs. Such a
representation need not reproduce the same sequential allocation on
the smaller state. The results here establish an elementary exact
criterion for a specified equilibrium family, not a replacement for
Miao's existence theorem.

\section{The declared restricted household problem}
Let $z,e$ have finite spaces, aggregate transition $q(z'\mid z)$, and
conditional individual transition $q(e'\mid e,z,z')$. Assets belong to
a Borel interval $A$. The reduced aggregate state is $s=(z,\mu)$,
where $\mu$ is a probability on $A\times E$ with finite first moment.
The capital and labor aggregates, prices and budget are those of the
catalogue incomplete-market model:
\begin{align*}
 K&=\int a\,d\mu>0, & L&=\int e\,d\mu>0,\\
 R&=1-\delta+F_K(z,K,L), & w&=F_L(z,K,L),\\
 c&=Ra+we-a'>0.
\end{align*}
Feasible savings satisfy the borrowing bound and any additional declared
asset constraint. Assume all relevant controls and kernels are Borel.
Impose exact conditional aggregation: for every feasible Borel savings rule
$g$, common state $s=(z,\mu)$, next aggregate shock $z'$, Borel set
$B\subseteq A$ and efficiency $e'$, the realized cross-sectional next
measure satisfies almost surely
\[
 \mu'(B\times\{e'\})=
 \sum_e\int_A \Ind\{g(a,e,s)\in B\}
 q(e'\mid e,z,z')\,\mu(da,\{e\}).
\]
This conditional continuum law of large numbers is an assumption on the
population probability construction, not a conclusion from individual
transition probabilities alone.
For the verification theorem impose a uniform bound $|u(c)|\leq M$
on every feasible consumption under every admissible deviation, and
$\beta\in(0,1)$. This holds, for example, on a declared compact
consumption interval bounded away from zero with continuous utility;
it is not automatic for logarithmic utility near zero.

\begin{theorem}[Recursive verification against history-dependent plans]
Let $D$ be an invariant Borel domain of aggregate states. Suppose Borel
prices, a Borel transition $\mathcal T(s,z')\in D$, a feasible savings
rule $g(a,e,s)$ and a bounded Borel function $V$ satisfy, at every feasible
individual state and every $s\in D$,
\[
 V(a,e,s)=\sup_{a'\in B(a,e,s)}
 \left\{u(R(s)a+w(s)e-a')+
 \beta\sum_{z',e'}q(z'\mid z)q(e'\mid e,z,z')
 V(a',e',\mathcal T(s,z'))\right\}, \tag{1}
\]
with $g$ attaining the supremum. Suppose $\mathcal T$ is the population
pushforward induced by $g$ and the efficiency kernel, and goods and
capital markets clear on $D$. Under the stated conditional continuum
law of large numbers, these objects generate a sequential competitive
equilibrium from every $s_0\in D$. The function $V$ is the actual
lifetime-utility supremum over all feasible adapted household plans,
not only Markov plans.
\end{theorem}
\begin{proof}
Recursively apply $q$ and $\mathcal T$ to construct the aggregate path;
for each household use the declared individual kernel and $g$. Their
Borel kernels determine finite-history laws consistently and hence the
path law. Distribution consistency and the continuum law of large
numbers give the specified cross-sectional distribution at each date;
the assumed pointwise clearing identities give market clearing.
For any household's adapted feasible plan, the aggregate path remains
fixed by price-taking. Conditional on its history, the continuation
kernel in (1) is the stated one. Bellman optimality gives
\[
 V(a_0,e_0,s_0)\geq
 \E\left[\sum_{t=0}^{T-1}\beta^t u(c_t)
                   +\beta^T V(a_T,e_T,s_T)\right].
\]
This follows by induction, multiplying each conditional inequality by
$\beta^t$ and cancelling consecutive continuation terms. Bounded $V$
makes the final term vanish uniformly. The utility bound makes the
discounted sum absolutely integrable and permits its limit inside
expectation. Thus $V$ bounds the utility of every adapted plan. For the
plan $g$, each Bellman inequality is equality; the same limit shows
its utility equals $V$. Feasibility, optimality and clearing are exactly
the sequential equilibrium conditions.
\end{proof}

\section{When a given sequential family factors through the state}
Take a nonempty countable set of initial aggregate states and, for each one, an
infinite-horizon sequential equilibrium specified at every positive-
probability finite aggregate history. Shocks have finite support, so the
union $\mathcal H$ of these histories is countable. Include continuations
after every positive-probability next shock, not just one realized path.
Let $s(h)=(z(h),\mu(h))$ and $D=\{s(h):h\in\mathcal H\}$.
Equip the probability-measure state space with its usual Borel structure;
its singletons are Borel, and thus every subset of countable $D$ is Borel.

The sequential equilibrium is assumed \emph{comprehensive}: at each
aggregate history $h$, a Borel current policy $g_h(a,e)$ is specified
and is part of a sequentially optimal continuation for every feasible
individual state, including states of zero current population mass.
Prices $(R_h,w_h)$ and aggregate next-state maps $s(hz')$ are also
specified. Individual state is sufficient for idiosyncratic transition
probabilities; the exogenous kernels above are independent of earlier
history conditional on their stated arguments. All household continuation
utilities obey the uniform bound used in the preceding theorem.

Define the \emph{collision conditions}: whenever $s(h)=s(\widetilde h)$,
\begin{align}
 (R_h,w_h)&=(R_{\widetilde h},w_{\widetilde h}), \tag{C1}\\
 g_h(a,e)&=g_{\widetilde h}(a,e)
       \quad\text{for every feasible }(a,e), \tag{C2}\\
 s(hz')&=s(\widetilde h z')
       \quad\text{for every }z'\text{ with }q(z'\mid z(h))>0. \tag{C3}
\end{align}
Equality on all feasible individual states is stronger than population-
almost-everywhere equality and is needed for the declared comprehensive
representation. Under the catalogue's factor-price formula (C1) is
automatic. Given distribution consistency, (C2) also implies (C3).

\begin{theorem}[Exact reduced-state representation on the reachable domain]
The specified comprehensive sequential family admits a single stationary
recursive representation $(R,w,g,V,\mathcal T)$ on $D$, agreeing with
its current policies, prices, aggregate transitions and continuation
utilities at every supported history and every feasible individual state,
if and only if (C1)--(C3) hold.
The resulting value satisfies $|V|\leq M/(1-\beta)$, equals the actual
sequential household optimum, and generates an equilibrium from every
initial state in $D$.
\end{theorem}
\begin{proof}
Necessity follows by evaluating each single-valued recursive function at
the same state in two histories. For sufficiency choose for each $s\in D$
one representative history $h_s$. Define prices and policy from this
history, and set $\mathcal T(s,z')=s(h_s z')$ for supported shocks.
For unsupported shocks choose any element of $D$; their values do not
affect expectations. Conditions (C1)--(C3) make all supported definitions
independent of the representative. Countability of $D$, together with
the Borel individual policies on each slice, proves measurability.
The domain is invariant by its definition from all supported finite
histories.

Starting from $(a,e,s)$, follow $g$ under these kernels and define $V$ to
be its expected discounted utility. Finite-horizon expected sums are
Borel by finite iteration of Borel kernels, and their tails are bounded
uniformly by $M\beta^T/(1-\beta)$. Their limit is therefore Borel and
bounded. The collision conditions imply by induction that the recursive
path and each original continuation have identical conditional laws:
their current policies agree, then (C3) agrees on next aggregate states,
and the individual transition kernel agrees by the primitive Markov
assumption. Thus this $V$ equals every corresponding sequential
continuation utility, including at zero-mass individual states.

For any feasible current $a'$, deviate once and then follow the assigned
optimal continuation at the next aggregate history and individual state.
Comprehensive sequential optimality makes the value of this deviation
at most $V(a,e,s)$. Its value is exactly the expression inside the
braces of (1), by the just-established continuation equality. Following
$g$ gives equality, by its definition as the achieved continuation
value. Hence (1) holds and $g$ attains it. Population consistency and
market clearing are inherited from the representatives and are
well-defined on collisions. The preceding verification theorem proves
sequential optimality and equilibrium from every $s\in D$.
\end{proof}

\begin{corollary}[What state enlargement must distinguish]
If two supported histories have the same reduced state but different
current policies at a feasible individual state, no comprehensive
recursive representation on $(z,\mu)$ can reproduce both continuations.
Any additional coordinate $v(h)$ that succeeds in reproducing them must
take different values at those histories. Conversely, if the map
$h\mapsto s(h)$ is injective, the collision conditions hold vacuously
(apart from primitive consistency), and the theorem gives a recursive
representation on its countable reachable domain.
\end{corollary}
\begin{proof}
A function evaluated at identical arguments cannot equal two different
policies. If an enlarged coordinate also agrees, its arguments remain
identical and the same contradiction applies. With injective $s$, each
state has one representative history, so all collision requirements are
automatic and the representation theorem applies.
\end{proof}

The last observation does not give a continuous or computationally useful
representation: a countable history tree may be encoded in a countable
set of distinct distributional states with very irregular neighboring
policies. It also does not give a common representation across all
uncountably many admissible initial distributions.

\section{Verification, relation to the target, and remaining gap}
The proof checks household deviations at zero-mass individual states;
checking only the equilibrium population path would not establish the
Bellman supremum. Boundedness is exactly what removes the terminal
continuation term; with unbounded utility it must be replaced by an
appropriate transversality or uniform-integrability condition for every
comparison plan. The collision obstruction is a logical criterion, not
an asserted example of multiplicity in the catalogue economy.
The first unresolved economic step is to prove collision consistency,
or construct its failure, from the specified incomplete-market primitives
without assuming a preselected comprehensive family. Extending the
measurable construction to an uncountable invariant domain requires a
measurable selection argument not supplied here. The full equivalence
and existence programme remains unresolved.

\begin{thebibliography}{9}
\bibitem{hsv} J. Heathcote, K. Storesletten and G. L. Violante (2009),
Quantitative Macroeconomics with Heterogeneous Households,
\emph{Annual Review of Economics} 1, 319--354, Section 5.2, footnote 20.
\url{https://doi.org/10.1146/annurev.economics.050708.142922}.
\bibitem{miao} J. Miao (2006), Competitive equilibria of economies with
a continuum of consumers and aggregate shocks,
\emph{Journal of Economic Theory} 128(1), 274--298.
\url{https://doi.org/10.1016/j.jet.2005.01.005}.
\end{thebibliography}

\end{document}
